% !TEX TS-program = lualatex

% FORCE-CROP

%%%
% src::
%   url = https://github.com/pfradin/luadraw/discussions/82#discussioncomment-14370152
%%%

\documentclass[border = 3pt]{standalone}

%%%
% Des couleurs faciles C'emploi via le package \xcolor!
%%%
\usepackage[svgnames]{xcolor}

%%%
% La bibliothèque \luadraw allie une facilité d’utilisation
% à un rendu particulièrement soigné.
%%%
\usepackage[3d]{luadraw}

\begin{document}

\begin{luadraw*}{name = circum-circle-sphere}
require 'luadraw_spherical'

local ld = luadraw

local Origin = ld.pt3d.Origin
local M      = ld.pt3d.M

local vecI = ld.pt3d.vecI
local vecJ = ld.pt3d.vecJ
local vecK = ld.pt3d.vecK

------
-- ¨Def de la zone graphique.
------
local graphview = ld.graph3d:new{
  window  = {-4, 4, -3, 4},
  viewdir = {30, 60},
  size    = {10, 10},
  bbox    = false
}

------
-- Réglages liés aux lignes et aux flèches.
------
graphview:Linewidth(4)

ld.Hiddenlines     = true
ld.Hiddenlinestyle = "dashed"

------
-- Nos ¨objs de ¨geo sphérique.
------
local O = Origin

local A, B, C = M(3, 0, 0), M(0, 3, 0), M(0, 0, 3)

local I1, R1, n1 = ld.circumcircle3d(A, B, O)
local I, R       = ld.circumsphere(O, A, B, C)

local T = ld.tetra(O, A - O, B - O, C - O)

graphview:Define_sphere({center = I, radius = R})

graphview:DScircle({I1, n1})
graphview:DSpolyline(ld.facetedges(T))

graphview:Dspherical()

graphview:Ddots3d({O, A, B, C, I, I1})

------
-- Ajout de points.
------
graphview:Dlabel3d(
  "$O$", O, {pos = "NW"},
  "$A$", A, {pos = "W"},
  "$B$", B, {pos = "E"},
  "$C$", C, {pos = "E"},
  "$I$", I, {pos = "E"},
  "$I_1$", I1, {pos = "S"}
)

------
-- Montrons le résultat de notre oeuvre.
------
graphview:Show()
\end{luadraw*}

\end{document}
